% =============================================================================
% DOCUMENTACIÓN DEL PAQUETE GLOSSARIES
% =============================================================================
% COMANDOS DISPONIBLES:
% \gls{RBD}      -> Primera vez muestra nombre completo + sigla; luego solo la sigla.
% \acrlong{RBD}  -> Siempre muestra la forma larga/expandida.
% \acrshort{RBD} -> Siempre muestra la forma corta/abreviada.
% \acrfull{RBD}  -> Siempre muestra la forma completa (larga + sigla).

% =============================================================================
% DEFINICIÓN DE ACRÓNIMOS DE LA TESIS
% =============================================================================
\newacronym{RBD}{RBD}{Diagrama de Bloques de Confiabilidad (\textit{Reliability Block Diagram})}
\newacronym{cRBD}{cRBD}{Diagrama de Bloques de Confiabilidad Coloreado (\textit{Colored Reliability Block Diagram})}
\newacronym{WMC}{WMC}{Conteo de Modelos Ponderado (\textit{Weighted Model Counting})}
\newacronym{CNF}{CNF}{Forma Normal Conjuntiva (\textit{Conjunctive Normal Form})}
\newacronym{WCNF}{WCNF}{Forma Normal Conjuntiva Ponderada (\textit{Weighted Conjunctive Normal Form})}
\newacronym{FT}{FT}{Árbol de Fallas (\textit{Fault Tree})}
\newacronym{FTs}{FTs}{Árboles de Falla (\textit{Fault Trees})}
\newacronym{DAG}{DAG}{Grafo Dirigido Acíclico (\textit{Directed Acyclic Graph})}